\documentclass[a4paper,12pt]{article}

% --- Packages ---
\usepackage[utf8]{inputenc}   % Handle UTF-8 characters
\usepackage{amsmath, amssymb} % Math symbols and environments
\usepackage{graphicx}         % Include images
\usepackage{hyperref}         % Clickable links
\usepackage{geometry}         % Adjust page margins
\geometry{margin=1in}

% --- Document info ---
\title{Vectorising the DHN-NCE loss function}
\author{Zain Nomani}
\date{\today}

% --- Begin document ---
\begin{document}

\maketitle

\begin{abstract}
This document summarizes my process of reconstructing and analyzing the loss function proposed in a healthcare AI research paper. 
The goal is to understand the underlying implementation, benchmark performance, and identify potential improvements.
\end{abstract}
\subsection{Loss Function}
The loss function was defined as:
\[
	\mathcal L = \mathcal L ^{v \to t} + \mathcal L ^{t \to v} \\
\]
\[
	\mathcal L ^{v \to t} = - \sum ^B _{i=1} \frac {I_{pi}^TT_{pi}}{\tau} + \sum ^B _{i=1} \log \sum _{j \ne i} \exp(\frac {I_{pi}^TT_{pj}}{\tau}) \mathcal W _{I_{pi}T_{pj}} ^{v \to t}
\]
\[
	\mathcal W = (B-1) \frac {\exp( \frac {\beta _1 I_{pi}^TT_{pj}}{\tau})} { \sum _{k \ne i} \exp(\frac {\beta _1 I_{pi}^TT_{pk}}{\tau})}
\]

And vice versa for $L^{t \to v}$

Since the weight matrix $W$'s leading diagonal ($i=j$) is not included in the second summation, we can manually set it to $-1$.
This will combine the 2 summations into 1, which we can begin to vectorise.

\[
	\mathcal L ^{v \to t} = \sum ^B _{i=1} \log \sum ^B _{j=1} \exp(\frac {I_{pi}^TT_{pj}}{\tau}) \mathcal W _{I_{pi}T_{pj}}^{v\to t}
\]
Where $\mathcal W _{ii} = -1 \space \forall i$. We can set this manually.

Let $S_{ij} = \frac {I_{pi}^TT_{pj}}{\tau} \quad$ where $I \in \mathbb R ^{B \times n}, \space T \in \mathbb R ^{B \times n}$ are the batch matrices of each image and text vector row-by-row.

So $S = \frac 1 \tau IT^T \in \mathbb R ^{B\times B}$ is a similarity matrix.

We normalise $I$ and $T$ to prevent gradient blow-up.

Next we define $A=\exp(S)$ element-wise.

So
$$
	\mathcal L^{v\to t} = \sum ^B _{i=1} \log \sum ^B _{j=1} A \mathcal W
$$
after setting $\text{diag}(\mathcal W) = -1$

Note that in order to calculate $\mathcal L ^{t\to v}$ we require $T_{pi}^TS_{pj}$, which is just $S^T$.

$$
	\mathcal W = (B-1) \frac {A_{ij}^{\beta _1}}{\sum _{k \ne i} A_{ij} ^{\beta _1}} = (B-1) \frac {e^{\beta _1 S_{ij}}}{\sum _{k \ne i} e^{\beta _1 S_{ik}}}
$$


We can define $X=\beta _1 S$ element-wise and set $X_{ii}=-\infty$

Then $\mathcal W^{v \to t} = (B-1) \cdot \text{softmax}(X, \dim = -1)$ computationally

We would transpose $X$ to find $\mathcal W^{t \to v}$

So our final vectorised algorithm looks like:

\begin{enumerate}

	\item Calculate $S = \frac 1 \tau IT^T$
 	\item Calculate $A = \exp (S)$ element-wise
 	\item Calculate $X = \beta _1 S$ and then set the leading diagonal to $-\infty$
	\item Calculate $\mathcal W = (B-1) \cdot \text{softmax} (X, \dim = -1)$
 	\item Set the leading diagonal of $\mathcal W$ to $-1$
 	\item $\mathcal L^{v\to t} = \sum (\log (\sum (A * \mathcal W)))$
 	\item Calculate $\mathcal L ^{t\to v}$ by repeating steps 4-6 with $X^T$ (to apply softmax to the other axis)

\[
	\mathcal L = \mathcal L^{t \to v} + \mathcal L^{v \to t}
\]
\end{enumerate}
\end{document}

